\chapter{Numerical Evaluation and Results Analysis}\label{ch:results}

This chapter tests the three links in the thesis contribution chain: whether
topology-aware recovery improves physically realizable sensing designs, whether
the robust formulation protects communication QoS under CSI uncertainty, and
whether the complete workflow remains useful after recovery and at a larger
network dimension. Each result is interpreted with its feasibility evidence,
not only with a conditional performance average.

\section{Evaluation Protocol, Metrics, and Baselines}

The experiments use the result sets stored with the paper figures. A run is
counted as physically feasible only after binary association recovery,
rank-one beamformer extraction or fixed-topology re-solution, and direct
validation of the recovered design. The validation checks the robust
communication SINR, sensing SINR, matrix PCRB, per-AP transmit-power limit,
and association cardinality. Consequently, an SDR solution or a continuous
association vector is not reported as a feasible physical design by itself.

All compared methods use common random scenarios within a sweep. Feasibility
rates are therefore unconditional statistics over those scenarios. Power,
runtime, rate, sensing-SINR, and PCRB summaries marked ``conditional'' are
computed only over the recovered designs that pass physical validation. The
baselines separate different design choices: FIM-greedy and nearest-AP methods
fix topology according to sensing geometry, random topology removes geometry
awareness, and relaxation or single-penalty variants isolate the recovery
mechanism. This protocol ensures that a low conditional power is never read as
an advantage when it is accompanied by frequent infeasibility.

\section{Topology-Recovery Advantage in Ordinary and Stress Geometries}

Figure~\ref{fig:method_nreq} compares topology-recovery strategies over 30
common seeds for a six-AP configuration. The proposed full recovery and the
FIM-greedy topology are physically feasible in all displayed cluster-size
settings. In contrast, the nearest-AP baseline is infeasible in 20\% of cases
at $N_{\mathrm{req}}=2$, and the random-topology baseline is substantially
less reliable for small and intermediate cluster sizes. The comparison shows
that the recovery workflow can identify physically usable sensing support where
purely local or random topology rules fail; the comparison is not evidence that
relaxation alone is exact.

The cardinality equality in (P1-C5') is a resource-allocation requirement, not
a monotone performance knob. Increasing $N_{\mathrm{req}}$ changes the
authorized sensing geometry and the feasible power allocation simultaneously,
so conditional power need not vary monotonically. The PCRB ratios remain close
to one in feasible recovered cases, which is consistent with the tracking
constraint being active at the solution.

\begin{figure}[htbp]
    \centering
    \includegraphics[width=\linewidth]{../../paper_results_and_figures_v1_0/figures/fig10_method_comparison_vs_nreq.png}
    \caption{Topology-method comparison versus required sensing APs per target.
    Feasibility is measured over all common scenarios; power, rate,
    sensing-SINR, and PCRB summaries are conditional on physically feasible
    recovered designs.}
    \label{fig:method_nreq}
\end{figure}

Figure~\ref{fig:pressure_geometries} asks whether the same topology advantage
persists when the geometry is deliberately unfavorable. With colocated edge
users and targets, and again with crowded targets, the proposed full recovery
and FIM-greedy topology remain feasible in all shown cases. The nearest-AP
topology loses feasibility in the edge-colocated setting and falls to roughly
24\% in the crowded setting. When it succeeds in the crowded case, its
conditional power approaches $300$~mW, whereas the proposed and FIM-greedy
methods remain near $50$~mW. The physical mechanism is geometric: selecting
the nearest APs alone can exhaust useful sensing directions or AP power under
crowding. The evidence is limited to the tested pressure configurations, but
it demonstrates why feasibility must accompany every conditional comparison.

\begin{figure}[htbp]
    \centering
    \includegraphics[width=\linewidth]{../../paper_results_and_figures_v1_0/figures/fig12_pressure_geometries.png}
    \caption{Physical feasibility and conditional power in two pressure
    geometries. Conditional-power bars exclude infeasible recovered runs and
    must be interpreted together with the associated feasibility rate.}
    \label{fig:pressure_geometries}
\end{figure}

\section{Robustness--Power Trade-Off under CSI Uncertainty}

Figure~\ref{fig:csi_robustness} compares the robust design with a nominal-CSI
baseline as the normalized CSI-uncertainty radius increases. Both methods are
certified feasible at their respective design points. When tested against
perturbations within the uncertainty ball, the robust design has zero observed
outage throughout the displayed sweep, whereas the nominal-CSI design has
approximately 90\% outage for nonzero uncertainty radii. The robust protection
raises transmit power moderately, from about $28.7$~mW at zero uncertainty to
about $30$~mW at the largest displayed radius.

This is the intended robustness--power trade-off: the additional power is
associated with markedly improved robustness under the tested perturbations.
It should not be interpreted as a guarantee outside the specified norm-bounded
uncertainty model, nor as a claim that the nominal design is infeasible at the
channel estimate itself.

\begin{figure}[htbp]
    \centering
    \includegraphics[width=\linewidth]{../../paper_results_and_figures_v1_0/figures/fig9_csi_robustness.png}
    \caption{Robustness comparison under increasing CSI-uncertainty radius.
    The middle panel is an empirical perturbation test and the right panel
    reports the corresponding transmit-power cost.}
    \label{fig:csi_robustness}
\end{figure}

\section{Recovery Mechanism and DC-SCA Stabilization}

The preceding comparisons concern recovered physical designs. Figure
\ref{fig:ablation} therefore isolates the mechanism that makes recovery
reliable. Across 25 common scenarios, all four variants have the same observed
physical-feasibility rate of 92\%, so the ablation does not show that one
penalty alone enlarges the feasible set in this test. Its informative result
is recovery quality: relaxation and rank-only variants retain a binary
distance on the order of $10^{-1}$, whereas binary-only and dual-DC variants
reduce the median binary distance below the $10^{-4}$ reference. The dual-DC
variant also maintains a very small median rank residual. Thus both penalties
are needed when the final design must be binary and rank one simultaneously.

\begin{figure}[htbp]
    \centering
    \includegraphics[width=\linewidth]{../../paper_results_and_figures_v1_0/figures/fig2_dual_dc_ablation.png}
    \caption{Recovery ablation over 25 common scenarios. Physical feasibility
    is evaluated after recovery; binary distance and rank residual diagnose
    whether a relaxed solution is suitable for physical recovery.}
    \label{fig:ablation}
\end{figure}

Figure~\ref{fig:dc_stabilization} provides the iteration-level complement to
the ablation on the same 25 seeds. The median binary distance falls from
approximately $0.43$ at initialization to below the recovery threshold after
the first DC-SCA iteration, and the recovery-ready fraction remains at 100\%
over the displayed iterations. The median rank residual stays at the
$10^{-9}$ scale. The continuous-power trace is a continuation diagnostic only;
it is not used to assert monotone descent of the original non-convex objective.

\begin{figure}[htbp]
    \centering
    \includegraphics[width=\linewidth]{../../paper_results_and_figures_v1_0/figures/fig5_statistical_double_dc_convergence.png}
    \caption{Statistical dual-DC SCA stabilization over 25 complete traces.
    The recovery-ready statistic combines a small binary distance, a small
    rank residual, and stable topology support.}
    \label{fig:dc_stabilization}
\end{figure}

\section{Larger-Scale Validation and Computational Cost}

Figure~\ref{fig:high_dimension} examines a larger configuration with $M=12$,
$K=6$, $P=3$, and $N_{\mathrm{req}}=3$. Over eight common seeds, the proposed
full recovery is physically feasible in all cases, whereas the FIM-greedy and
nearest-AP baselines achieve 75\%, and the random topology is feasible in only
12.5\% of cases. The proposed conditional energy cost is close to the
FIM-greedy cost and substantially below the nearest-AP and random topologies.
Its wall-clock cost is higher than that of fixed-topology baselines because the
workflow performs candidate generation, fixed-topology re-solves, and physical
validation.

These observations provide scalability evidence at the tested dimensions, not
an asymptotic complexity claim. The smaller number of seeds also means that
larger independent sweeps and confidence intervals are required before drawing
stronger quantitative conclusions.

\begin{figure}[htbp]
    \centering
    \includegraphics[width=\linewidth]{../../paper_results_and_figures_v1_0/figures/fig11_m12_scalability_validation.png}
    \caption{Scalability validation for $M=12$, $K=6$, $P=3$, and
    $N_{\mathrm{req}}=3$ over eight common seeds. Power and runtime are
    conditional summaries for physically feasible recovered designs.}
    \label{fig:high_dimension}
\end{figure}

\section{Summary of Findings}

The numerical evidence supports three bounded conclusions. First,
topology-aware recovery improves physical feasibility relative to nearest and
random topology rules in the tested ordinary and pressure geometries. Second,
the robust formulation exchanges a modest power increase for substantially
improved outage behaviour within the stated CSI-uncertainty model. Third, the
dual DC penalties provide recovery-quality evidence, while the larger-scale
experiment indicates that the validated workflow remains usable beyond the
smallest configuration. These conclusions remain empirical: they do not prove
global optimality, universal relaxation tightness, or asymptotic scalability.
